% ============================================================
%  02_kien_thuc_nen_tang.tex  --  Kien thuc nen tang
% ============================================================

\section{Kiến thức nền tảng}
\label{sec:kienthuc}

\subsection{Các lớp lỗi quản lý bộ nhớ trong C}

Quản lý bộ nhớ thủ công trong C là nguồn gốc của nhiều lỗ hổng bảo mật nghiêm
trọng nhất trong lịch sử phần mềm. Ba lớp lỗi được nghiên cứu trong bài thực
hành này đều thuộc danh sách CWE Top 25 Most Dangerous Software Weaknesses.

\subsubsection{Buffer Overflow -- CWE-121}

Buffer Overflow (tràn bộ đệm) xảy ra khi chương trình ghi dữ liệu vượt ra ngoài
ranh giới của vùng nhớ đã cấp phát. Trong C, nguy cơ này xuất hiện phổ biến
nhất khi dùng các hàm sao chép chuỗi không kiểm tra độ dài như \textit{strcpy()},
\textit{gets()}, \textit{sprintf()} với dữ liệu đầu vào không đáng tin cậy.

\begin{itemize}
    \item Cơ chế: Vùng đệm được cấp phát $N$ byte; dữ liệu ghi vào có
          độ dài $M > N$. Phần dữ liệu dư ($M - N$ byte) ghi đè lên vùng nhớ
          liền kề, có thể là metadata heap, địa chỉ trả về trên stack, hoặc dữ
          liệu của biến khác.
    \item Hậu quả: Trong trường hợp nhẹ: dữ liệu bị hỏng, chương trình
          crash. Trong trường hợp nghiêm trọng: kẻ tấn công có thể ghi đè địa
          chỉ trả về để điều khiển luồng thực thi (Return-Oriented Programming),
          dẫn đến thực thi mã tùy ý (Arbitrary Code Execution).
    \item Trong module Inventory: Hàm \textit{add\_book\_title\_UNSAFE\_HEAP()}
          dùng \textit{strcpy(buf, title)} với \textit{buf} chỉ có 32 byte. Payload
          dài hơn 31 ký tự sẽ ghi vượt ranh giới heap buffer.
\end{itemize}

\subsubsection{Use-After-Free -- CWE-416}

Use-After-Free (sử dụng sau khi giải phóng) xảy ra khi chương trình tiếp tục
đọc hoặc ghi thông qua một con trỏ trỏ đến vùng nhớ đã được \textit{free()}.

\begin{itemize}
    \item Cơ chế: Sau khi \textit{free(ptr)}, bộ quản lý heap có thể
          tái sử dụng vùng nhớ đó cho một allocation khác. Nếu chương trình vẫn
          đọc/ghi qua \textit{ptr}, nó thực ra đang thao tác trên dữ liệu của
          đối tượng mới, gây ra hành vi không xác định (Undefined Behavior).
    \item Hậu quả: Kẻ tấn công có thể kiểm soát nội dung vùng nhớ được
          cấp phát lại (heap grooming/heap spray), sau đó kích hoạt đọc/ghi qua
          con trỏ cũ để thực thi mã tùy ý hoặc leo thang đặc quyền.
    \item Trong module Inventory: Hàm \textit{demo\_use\_after\_free()}
          gọi \textit{free(b)} rồi ngay sau đó đọc \textit{b->year}. ASan phát
          hiện đây là \textit{heap-use-after-free}.
\end{itemize}

\subsubsection{Memory Leak -- CWE-401}

Memory Leak (rò rỉ bộ nhớ) xảy ra khi bộ nhớ được cấp phát động nhưng không
bao giờ được giải phóng, ngay cả khi không còn tham chiếu nào đến vùng nhớ đó.

\begin{itemize}
    \item Cơ chế: Mỗi lần vòng lặp chạy, \textit{malloc()} cấp phát
          một vùng nhớ mới; con trỏ bị ghi đè bởi vòng lặp tiếp theo mà không
          gọi \textit{free()} trước. Vùng nhớ cũ mất đi tham chiếu, không thể
          giải phóng và tích lũy liên tục.
    \item Hậu quả: Trong ứng dụng chạy ngắn, tác động thường không đáng
          kể. Trong hệ thống phục vụ lâu dài (daemon, server), rò rỉ tích lũy
          dẫn đến cạn kiệt bộ nhớ và Denial of Service.
    \item Trong module Inventory: Hàm \textit{demo\_memory\_leak(count)}
          cấp phát \textit{count} đối tượng \textit{Book} trong vòng lặp nhưng
          không gọi \textit{free()} ở bất kỳ đâu.
\end{itemize}

\subsubsection{NULL Pointer Dereference -- CWE-476 (tiềm ẩn)}

Khi \textit{malloc()} thất bại (hệ thống hết bộ nhớ), nó trả về \textit{NULL}.
Nếu chương trình không kiểm tra và tiếp tục dùng con trỏ NULL để gọi
\textit{strcpy(NULL, title)}, kết quả là crash ngay lập tức -- một dạng
Denial of Service đơn giản.

% ────────────────────────────────────────────────────────────
\subsection{AddressSanitizer}

AddressSanitizer (ASan) là công cụ phân tích bộ nhớ động được phát triển bởi
Google, tích hợp trong GCC từ phiên bản 4.8 và Clang từ 3.1. Khi biên dịch với
cờ \textit{-fsanitize=address}, trình biên dịch chèn thêm mã kiểm tra vào mỗi
thao tác đọc/ghi bộ nhớ và mỗi lời gọi \textit{malloc}/\textit{free}.

\subsubsection{Cơ chế hoạt động}

\begin{itemize}
    \item Shadow memory: ASan duy trì một vùng bộ nhớ bóng (shadow
          memory) ánh xạ trạng thái của từng byte: có thể truy cập an toàn,
          là redzone (vùng bảo vệ) hay đã được giải phóng.
    \item Redzone: Quanh mỗi vùng cấp phát (cả stack và heap), ASan
          chèn các vùng đệm bảo vệ được đánh dấu trong shadow memory. Mọi
          truy cập vào redzone đều sinh lỗi ngay lập tức.
    \item Quarantine: Sau khi \textit{free()}, vùng nhớ được đặt vào
          danh sách cách ly và đánh dấu ``đã giải phóng'' trong shadow memory
          thay vì trả về cho allocator ngay. Mọi truy cập vào vùng này sinh
          lỗi \textit{heap-use-after-free}.
    \item LeakSanitizer: Khi bật \textit{ASAN\_OPTIONS=detect\_leaks=1},
          LeakSanitizer quét toàn bộ bộ nhớ khi chương trình kết thúc, tìm
          các vùng đã cấp phát nhưng không còn tham chiếu nào trỏ đến.
\end{itemize}

\subsubsection{Giới hạn quan trọng: tràn trong struct}

ASan đặt redzone ở \emph{biên} của mỗi vùng cấp phát -- không đặt
\emph{giữa} các trường của cùng một allocation. Xét struct sau:

\begin{lstlisting}[language=C, caption={Struct Book -- hai trường liền kề trong một allocation}]
typedef struct {
    char title[32];   /* offset 0..31  */
    char author[32];  /* offset 32..63 */
    int  year;
    int  copies_available;
} Book;
\end{lstlisting}

Khi khai báo \textit{Book b} trên stack, toàn bộ struct là một allocation duy
nhất. Redzone được đặt \emph{trước} \textit{b.title} và \emph{sau}
\textit{b.copies\_available}. Nếu \textit{strcpy(b.title, payload)} ghi 35 ký
tự, 3 byte dư ghi vào \textit{b.author[0..2]} -- vẫn trong vùng allocation,
không đụng đến redzone -- nên ASan không phát hiện.

Ngược lại, khi cấp phát \textit{buf = malloc(32)} riêng biệt, redzone nằm ngay
sau byte thứ 31 của \textit{buf}. Ghi byte thứ 32 vào đó ngay lập tức đụng
redzone và ASan dừng chương trình với thông báo \textit{heap-buffer-overflow}.

Kết luận sư phạm: Không có lỗi ASan $\neq$ chương trình an toàn.
Kiểm tra không đầy đủ (thiếu LeakSanitizer, thiếu kịch bản edge-case) hoặc bố
trí bộ nhớ nhất định có thể khiến lỗi thật sự tồn tại mà không bị phát hiện.

% ────────────────────────────────────────────────────────────
\subsection{Nguyên tắc lập trình C an toàn}

Để phòng tránh các lỗi bộ nhớ, cần tuân thủ các nguyên tắc sau:

\begin{itemize}
    \item Không bao giờ dùng \textit{strcpy} với đầu vào không tin cậy.
          Thay bằng \textit{snprintf(dst, size, "\%s", src)} -- hàm này luôn giới
          hạn số byte ghi và đảm bảo byte kết thúc chuỗi \textit{'\textbackslash0'}.
    \item Kiểm tra giá trị trả về của \textit{malloc}. Nếu \textit{malloc}
          trả về \textit{NULL}, dừng thực thi hoặc xử lý lỗi trước khi dùng con
          trỏ.
    \item Đặt con trỏ về \textit{NULL} ngay sau \textit{free}. Mọi truy
          cập tiếp theo qua con trỏ NULL sẽ bị phát hiện ngay (NULL dereference
          crash), thay vì tạo ra hành vi không xác định nguy hiểm hơn.
    \item Mỗi \textit{malloc} phải có một \textit{free} tương ứng. Cần
          phân tích rõ ownership: ai cấp phát vùng nhớ, ai chịu trách nhiệm giải
          phóng, và trong trường hợp lỗi có đường thoát sớm nào bị bỏ sót không.
    \item Sao chép dữ liệu cần dùng trước khi giải phóng. Nếu cần đọc
          giá trị từ struct sau này, sao chép vào biến cục bộ trước khi gọi
          \textit{free()}.
\end{itemize}

% ────────────────────────────────────────────────────────────
\subsection{Mô hình đánh giá rủi ro CVSS-lite}

Bài thực hành áp dụng một mô hình đơn giản hóa lấy cảm hứng từ CVSS v3.1
(Common Vulnerability Scoring System). Đây là mô hình giáo dục, \emph{không}
phải công thức CVSS chính thức.

\subsubsection{Các trục đánh giá}

\begin{itemize}
    \item Attack Vector (AV): Mức độ tiếp cận cần thiết để khai thác
          -- NETWORK (0.85), ADJACENT (0.62), LOCAL (0.55), PHYSICAL (0.20).
    \item Attack Complexity (AC): Độ phức tạp của cuộc tấn công --
          LOW (0.77), HIGH (0.44).
    \item Privileges Required (PR): Quyền hạn cần có -- NONE (0.85),
          LOW (0.62), HIGH (0.27).
    \item User Interaction (UI): Có cần tương tác người dùng không --
          NONE (0.85), REQUIRED (0.62).
    \item Confidentiality / Integrity / Availability (C/I/A):
          Mức độ ảnh hưởng đến từng trục -- NONE (0.00), LOW (0.22), HIGH (0.56).
\end{itemize}

\subsubsection{Công thức tính điểm}

\[
    \mathrm{Exploitability}_{sub} = 8.22 \times AV \times AC \times PR \times UI
\]
\[
    \mathrm{ISC}_{base} = 1 - (1-C)(1-I)(1-A)
\]
\[
    \mathrm{Impact}_{sub} = 6.42 \times \mathrm{ISC}_{base}
\]
\[
    \mathrm{BaseScore} = \min(\mathrm{Impact}_{sub} + \mathrm{Exploitability}_{sub},\ 10.0)
\]

Điểm BaseScore được xếp hạng mức độ theo thang CVSS chuẩn: NONE (0.0), LOW
(0.1--3.9), MEDIUM (4.0--6.9), HIGH (7.0--8.9), CRITICAL (9.0--10.0).
